% CHECKPOINT33: new source; no checkpoint32 file is changed.
\statement{A.104}{Proved}{The actual congruence automorphism and the exact kernel-surjectivity gate}
Retain the actual \(u,k,d,q,U,b\) of A.103 and put
\[
 \begin{gathered}
 \psi_d:F\longrightarrow F,\qquad
 \psi_d(x)=x^d,\quad \psi_d(y)=u^{-1}y^du,\\
 C=q(F),\qquad N=\ker q,\qquad
 H=\begin{pmatrix}1&-b\\0&d^{-1}\end{pmatrix}.
 \end{gathered}
 \tag{A.104.1}
\]
Then \(H\) normalizes the actual closed congruence image \(C\), and
\[
 q\psi_d=\operatorname{Ad}(H)q,\qquad
 \psi_d^{-1}(N)=N.
 \tag{A.104.2}
\]
Writing \(L=\psi_d(F)\), one has the exact equalities
\[
 \begin{gathered}
 L=\overline{\langle x,u^{-1}yu\rangle},\qquad
 F=NL,\qquad L\cap N=\psi_d(N),\\
 L=F\quad\Longleftrightarrow\quad \psi_d(N)=N.
 \end{gathered}
 \tag{A.104.3}
\]
Thus the remaining full-surjectivity question is exactly
surjectivity of the actual restriction \(\psi_d:N\to N\).
No surjectivity on that infinite kernel is asserted.

For actual normalized full-GT data
\[
 \alpha(x)=x^{-1},\quad
 \alpha(y)=\varphi^{-1}y^{-1}\varphi,\quad
 u=\varphi\alpha(u),
 \tag{A.104.4}
\]
let \(c_0\) denote the automorphism inverting both \(x,y\), and put
\(J_0=\operatorname{diag}(1,-1)\).
The original involution preserves \(N\), and its actual action on
\(C\) is
\[
 \begin{gathered}
 q\alpha=\operatorname{Ad}(K)q,\qquad
 K=HJ_0H^{-1}=\begin{pmatrix}1&2bd\\0&-1\end{pmatrix},\\
 q(\varphi)=Y^{-bd}X^{bd}
       =\begin{pmatrix}1&2bd\\2bd&1+4b^2d^2\end{pmatrix}.
 \end{gathered}
 \tag{A.104.5}
\]
In particular the quotient involutions are conjugate by
\(\operatorname{Ad}(H)\). This conclusion is proved directly from
the original peripheral action and norm, without cancelling an
unproved automorphism \(\psi_d\).

\begin{proof}
Because \(d\) is a unit, \(H\) is invertible with
\[
 H^{-1}=\begin{pmatrix}1&bd\\0&d\end{pmatrix}.
 \tag{A.104.6}
\]
Exact multiplication over \(\widehat{\mathbb Z}\) gives
\[
 \begin{gathered}
 HXH^{-1}=X^d,\\
 HYH^{-1}
   =\begin{pmatrix}
        1+2b&2db^2\\-2d^{-1}&1-2b
     \end{pmatrix}
   =U^{-1}Y^dU.
 \end{gathered}
 \tag{A.104.7}
\]
The last equality uses the full actual matrix of A.103.
Equality on the two free generators proves (A.104.2)'s first
identity on the entire profinite group.

For every positive integer \(m\), let
\(C_m=\langle X_m,Y_m\rangle=q_m(F)\), the actual finite image.
Equation (A.104.7) implies
\(H_m C_mH_m^{-1}\subseteq C_m\).
Conjugation is injective and \(C_m\) is finite, so this inclusion
is equality. The compact subgroup \(HCH^{-1}\) therefore has
the same projection as \(C\) in every congruence quotient.
Closed subgroups of this inverse limit are determined by those
projections, proving \(HCH^{-1}=C\).
No classification of \(C\), and no assumption that it is the
entire ambient \(\operatorname{SL}_2\), is needed.
Conjugation by \(H\) is an external normalizer action: \(H\) is
not asserted to lie in \(C\).

For any \(g\in F\), the first identity in (A.104.2) gives
\[
 \psi_d(g)\in N
 \quad\Longleftrightarrow\quad
 Hq(g)H^{-1}=I
 \quad\Longleftrightarrow\quad g\in N.
 \tag{A.104.8}
\]
The image \(L\) is compact, hence closed.
Unit powering does not change either generated procyclic subgroup,
so its description in (A.104.3) is literal.
Since \(q(L)=C=q(F)\), every \(g\in F\) has the same \(q\)-image
as some \(\ell\in L\); then \(g\ell^{-1}\in N\).
This proves \(F=NL\).
Also
\[
 L\cap N
   =\psi_d\bigl(\psi_d^{-1}(N)\bigr)
   =\psi_d(N).
 \tag{A.104.9}
\]
If \(\psi_d(N)=N\), then \(N\subseteq L\) and \(F=NL=L\).
Conversely \(L=F\) makes (A.104.9) equal to \(N\).
This proves the exact criterion without any finite-generation
assumption on \(N\).

\emph{The original involution, with no cancellation through \(L\).}
Let
\(\Phi=q(\varphi)\).
The defining GT0 two-cycle relation for \(\varphi\), evaluated
using \(S_0\) from A.103, makes \(\Phi\) symmetric. Write
\[
 \Phi=\begin{pmatrix}a&t\\t&e\end{pmatrix}.
 \tag{A.104.10}
\]
Original third-peripheral inversion means that \(\alpha(xy)\)
is conjugate to \((xy)^{-1}\) in \(F\).
Evaluate this statement directly under \(q\), without presuming
that \(q\)'s kernel is invariant. Taking traces gives
\[
 \operatorname{tr}(X^{-1}\Phi^{-1}Y^{-1}\Phi)
      =2-4a^2
      =\operatorname{tr}((XY)^{-1})=-2.
 \tag{A.104.11}
\]
The additive group \(\widehat{\mathbb Z}\) is torsion-free, so
\(a^2=1\). Thus \(a\) is a unit and
\(e=a(1+t^2)\).
Set
\[
 K_\Phi=\begin{pmatrix}1&at\\0&-1\end{pmatrix},
 \qquad K_\Phi^2=I.
 \tag{A.104.12}
\]
Directly on the original two generators,
\[
 \begin{gathered}
 K_\Phi X K_\Phi^{-1}=X^{-1},\\
 K_\Phi Y K_\Phi^{-1}
   =\begin{pmatrix}1-2at&-2t^2\\2&1+2at\end{pmatrix}
   =\Phi^{-1}Y^{-1}\Phi.
 \end{gathered}
 \tag{A.104.13}
\]
Consequently \(q\alpha=\operatorname{Ad}(K_\Phi)q\) on all
of \(F\), which proves \(N\)-invariance directly.
This intermediate descent uses only the actual negative
peripheral action and the original two-cycle relation.

Now use the actual canonical norm.
Its image is \(U=\Phi K_\Phi U K_\Phi\), or
\[
 UK_\Phi=\Phi K_\Phi U,\qquad
 \Phi K_\Phi=\begin{pmatrix}a&0\\t&-a\end{pmatrix}.
 \tag{A.104.14}
\]
Since \(U_{11}=d^{-1}\) is a unit, comparison of the
\((1,1)\)-entries gives \(a=1\).
The \((1,2)\)-entries then give
\(d^{-1}t-b=b\), hence \(t=2bd\).
Substitution proves every formula in (A.104.5), including
\(K_\Phi=HJ_0H^{-1}\).
The standard inversion satisfies
\(qc_0=\operatorname{Ad}(J_0)q\); therefore the induced
involutions on the actual quotient are indeed conjugate.

\emph{Comparison with the other quotient gates.}
Put \(\psi_u(x)=x,\ \psi_u(y)=u^{-1}yu\).
The power substitution \(p_d:x\mapsto x^d,\ y\mapsto y^d\)
is a genuine automorphism, with inverse \(p_{d^{-1}}\), and
\(\psi_d=\psi_u p_d\).
Thus the normalized map has exactly the original canonical
surjectivity question.
Its abelianization is \(dI_2\), so on each maximal free pro-\(p\)
quotient it is an automorphism by the Frattini criterion.
That conclusion and the congruence automorphism have different
quotient scopes: neither implication is being inferred from the
other. In particular nothing here proves automorphy on a relative
pro-\(p\) extension above a nontrivial congruence head.

The actual norm also gives
\(\alpha\psi_d=\psi_dc_0\), directly on \(x,y\).
The proof of (A.104.5) did not cancel \(\psi_d\) in this relation.
Any properness of \(L\) is invisible to all the signed congruence
heads, because \(q_m(L)=q_m(F)\) for every \(m\).
A.103 places the original marked defect \(H_d\) in \(N\);
it does not evaluate that word inside \(N\).
Neither the kernel-surjectivity criterion nor the quotient
conjugacy proves the full original cusp equation or a full
automorphism of \(F\).
\end{proof}
